
\begin{proposicion}[A distant minor; stability \textup{(S)}]\label{prop:menor-lejano}
For constants $(C_p,\gamma_A)$ there are a threshold $k_S=k_S(C_p,\gamma_A)\ge k_2$ and a constant
$C_{S1}=C_{S1}(C_p,\gamma_A)$ with the following properties.
\begin{enumerate}
\item[\textup{(a)}] Let $N$ satisfy \textup{(H)} with $(C_p,\gamma_A)$ and have size $k\ge k_S$, let $d$
be a cell of $N$ and $c$ a cell of $N_d$. Then $|r_c(N_d)-r_c(N)|\le C_{S1}/k^3$.
\item[\textup{(b)}] Let $M$ satisfy \textup{(H)} with $(C_p,\gamma_A)$ and have size $k$, with
$k\ge k_S(C_p,\gamma_A)$ and $k-1\ge k_S(\tfrac54C_p,\tfrac34\gamma_A)$. Let $e,f$ lie in distinct rows
and columns and let $c$ lie in none of their rows or columns. Then
\[
 |r_c(M_{ef})-r_c(M)|\le\frac{C_S}{k^3}, \tag{S}
\]
with $C_S$ depending only on $(C_p,\gamma_A)$.
\end{enumerate}
In the class $\mathcal K$ both hold with $(C_p,\gamma_A)=(\tfrac32M_0,0.88\gamma)$.
\end{proposicion}
\begin{proof}
Part (b) follows from (a). Indeed $M_{ef}=(M_e)_f$; by Step~4 of the proof of
Theorem~\ref{teo:insercion-torres}, $M_e$ satisfies \textup{(H)} with $(\tfrac54C_p,\tfrac34\gamma_A)$ at
size $k-1$; and $c$ is a cell of $M_e$ and of $M_{ef}$. Apply (a) to $(M,e)$ and to $(M_e,f)$ and add.

For (a) we use the notation of that proof with $N$ in place of $M$ and $e_0=c$, write $h=1/k$, and let
$\delta$ denote the variation of a quantity from $N$ to $N_d$, read on the common support. Constants
depend only on $(C_p,\gamma_A)$, and $k_S$ is chosen so that the elementary inequalities used below hold.

\emph{Step 1. A second-order representative of $r_{e_0}$.} Let $u_x=a_x^\top\xi_N$, $\alpha_x=\mathsf
G(x,x)-2$, $\eta_x=\mathsf G(x,e_0)$, $b_x=\tfrac12\tau_x\dot\tau_xu_x^2$ with $\dot\tau_x=\dot\tau(p_x)$,
and
\[
 Z=(S'(\tau|)+\Pi')^{-1}\Bigl(\rho-\sum_{x\in M'}b_xa_x\Bigr),\qquad z_x=a_x^\top Z .
\]
These are computed from the data of the parent. The mass equation of the child is exactly
\begin{equation}\label{eq:masas-hijo}
 \zeta=(S'(\tau|)+\Pi')^{-1}(\rho-B),\qquad
 B_i=\sum_{x\in i\cap M'}\bigl[u(t_xe^{s_x})-p_x-\tau_xs_x\bigr].
\end{equation}
Since $\|\xi\|_\infty=O(h)$ and $\|\xi-\xi_N\|_\infty=O(h^2)$ by~\eqref{eq:newton-hijo}, the third
derivative of $s\mapsto u(te^s)$ is $O(p)$ on the segment, the line masses are bounded and the inverse is
bounded in the supremum norm, a Taylor expansion gives $\zeta=Z+O_\infty(h^3)$.

In~\eqref{eq:re-restos} we replace the remainders as follows; $E_{\mathcal D3}=\zeta^\top\rho$ and
$E_{\mathcal D4}=-\tfrac12\zeta^\top S'(\tau|)\zeta$ are already $O(h^3)$.
\begin{enumerate}
\item In $E_{\rm main}$, replace $p_x^3/(1-2p_x)$ by $p_x^3$.
\item Replace $E_{\mathcal D1}=\tfrac12\sum_{F^*}\varsigma_x\beta_{\lambda(x)}$ by
$\tfrac12\sum_{F^*}p_x^2\beta_{\lambda(x)}$.
\item In $E_{\mathcal D2}=-\tfrac12\beta_I^\top\rho$, replace $\rho_{\lambda(x)}$ by
$-p_x^2-p_x(\eta_x-1)$ for $x\in F^*$.
\item Replace $E_{\mathcal D5}=-R_3$ by $-\tfrac16\sum_{M'}\tau_x\dot\tau_xu_x^3$.
\item Replace $E_{L1}$ by $-\tfrac23\sum_{F^*}p_x^3-p_0^2-\tfrac14(\mu_r^2+\mu_c^2)$.
\item Replace $E_{L2}$ by $\sum_{M'}\bigl\{\tfrac12D(\ell_2'\tau)_xu_x^2
+(\ell_2'(p_x)-p_x)\tau_xu_x+p_x\tau_xz_x\bigr\}$, where $\ell_2'(u)=u/[(1-2u)(1-u)]$ and
$D=t\,d/dt$, so that $Du=\tau$.
\item Replace $E_{\rm(i)}$ by $-\sum_{M'}\bigl\{\tau\alpha z+4p\tau z+b\alpha
+2(\dot\varsigma-2p)\tau u+2\dot\varsigma b+\ddot\varsigma(\tau u)^2+\dot\varsigma\tau u\alpha\bigr\}_x$.
\item Replace $E_{L3}$ by $\tfrac14\operatorname{tr}(\mathsf S^+W_0\mathsf S^+W_0)$, where $W_0$ is
obtained from the entries of $W$ in Step~7 by setting $\Delta\varsigma_x=\dot\varsigma_x\tau_xu_x$ on the
diagonal and $\tau'_x-\tau_x=\dot\tau_x\tau_xu_x$ off the diagonal.
\item Replace $E_{\delta1}$ by $\tfrac12\hat c_{e_0}^2$.
\item Replace $E_{\delta2}=p_0\hat c_{e_0}/(1-p_0)$ by $p_0\hat c_{e_0}$.
\item In
$E_X=\sum_{M'}\tau_x\tilde\varepsilon_x\eta_x+\sum_{F^*}\tau_x\tilde\varepsilon_x
+\sum_{F^*}\tau_x\varepsilon_x(\eta_x-1)+\tau_0\varepsilon_0\mathsf G_0-\tilde\varepsilon_0$, replace
$\tilde\varepsilon_x$ by $6p_x^2+2p_x\alpha_x$ and $\varepsilon_x$ by $\alpha_x+2p_x$.
\end{enumerate}
Let $r^{[2]}_{e_0}$ be the sum of these expressions with the coefficients of~\eqref{eq:re-restos}. Each
replacement costs $O(h^3)$, so $r_{e_0}=r^{[2]}_{e_0}+O(h^3)$ uniformly, in $N$ and in $N_d$. To see this,
note that on $F^*$ we have $\sum p\le2$ and $\sum p^j=O(h^{j-1})$, and on $M'$ we have $\sum p\le k$,
$\sum p^2=O(1)$, $\sum p^3=O(h)$. The errors in items 1--3 are $O(\sum_{F^*}p^4)$ and
$O(h\sum_{F^*}(p^3+hp^2))$, and that of item~4 is $O(h^4\sum_{M'}p)$. For item~5,
$-\log(1-u)-u+\tfrac12c_2(u)=-\tfrac23u^3+O(u^4)$; the terms $\tfrac12c_2(p_0)$ and
$\tfrac12\log(1-\tau_0^2/(m_rm_c))$ contribute $-\tfrac12p_0^2+O(h^3)$ each, and the logarithms of the two
masses give $-\tfrac14(\mu_r^2+\mu_c^2)+O(h^3)$. For item~6, $D(\ell_2'\tau)$ and $D^2(\ell_2'\tau)$ are
$O(p^2)$ and $(\ell_2'-p)\tau=O(p^3)$. For item~7, with $s=a_x^\top\xi$,
\[
 \Delta p-\tau s-b=O(ph^3),\qquad \Delta\varsigma-\dot\varsigma\tau u=O(p^2h^2),\qquad
 \Delta\varsigma-\dot\varsigma\Delta p-\tfrac12\ddot\varsigma(\tau u)^2=O(p^2h^3),
\]
and one sums using $\dot\varsigma=O(p)$, $\dot\varsigma-2p=O(p^2)$, $\alpha=O(h)$ and $|s-u|=O(h^2)$. For
item~8 the error on the diagonal is $O(h^2\sum_{x\in i}p_x^2)=O(h^3)$ per line and $O(ph^2)$ per entry off
the diagonal, so $\|W-W_0\|_{\rm HS}=O(h^2)$; since $\|W\|_{\rm HS}=O(h)$, the quadratic trace changes by
$O(h^3)$, and the tail of the logarithm is $O(h^3)$ by the spectral bound. Items 9 and 10 follow from
$-\log(1+x)+x=\tfrac12x^2+O(x^3)$ and $\hat c_{e_0}=O(h)$. For item~11,
$\dot c_2+\dot\tau(2+\alpha)-2=\alpha+2p+6p^2+2p\alpha+O(p^3+p^2|\alpha|)$, and the sums are controlled
by $\sum_{M'}\tau|\eta|\le C$, $|\eta-1|\le Ch$ on $F^*$ and $\tau_0|\mathsf G_0|=O(h)$.

\emph{Step 2. Variation of the displacement.} Theorem~\ref{teo:insercion-torres} and
Corollary~\ref{cor:dos-casillas}, applied to the minor of $d$, give on the common support
\begin{equation}\label{eq:variaciones-base}
 \delta p=O(ph),\qquad \delta G^0_{ij}=O(h^2),\qquad \delta\hat c=O(h^2),\qquad
 \delta\tau=O(ph),\qquad \delta\mu_i=O(h^2),
\end{equation}
the last by summation over a line that loses at most one site. From $\zeta=O(h^2)$ and
$s_x=-\mathsf G(x,e_0)+a_x^\top\zeta$ we first get $\delta s=O(h^2)$. In~\eqref{eq:masas-hijo} the summand
of $B_i$ is $p_xs_x^2$ times a function with bounded derivatives in $(p_x,s_x)$ near the origin. On the
common indices it changes by $O(p_xh^3)$, and the lost site contributes $O(h^3)$; thus $\delta
B_i=O(h^3)$. On the star $F^*$, the exact formula~\eqref{eq:rho-hijo} gives $\delta\rho_i=O(h^3)$. Of the
two linear systems~\eqref{eq:masas-hijo}, for $N$ and for $N_d$, restrict the larger one to the common
indices. The absolute row sums of the difference of the two matrices, including the coupling with the two
removed indices, are $O(h)$: this follows from $\delta\tau=O(ph)$, from the lost site, and from the change
of the gauge projection. Multiplying by the solution, which is $O(h^2)$, and by the inverse, which is
uniformly bounded in the supremum norm, we obtain
\[
 \delta\zeta=O(h^3),\qquad\text{and in the same way}\qquad \delta z=O(h^3),\qquad \delta u=O(h^2).
\]

\emph{Step 3. Variation of the remainder.} By Step~1 it suffices to compare the eleven expressions in $N$
and in $N_d$. The interior terms of items 4, 6 and 7 have majorants $Cph^3$, $Cp^2h^2$ and $Cp^3h$ per
site; each sums to $O(h^2)$, and to $O(h^3)$ over the two lost lines. On the common indices the
telescoping of each product gains a factor $h$ in the majorant, by~\eqref{eq:variaciones-base},
$\delta\alpha=O(h^2)$, $\delta z=O(h^3)$ and $\delta u=O(h^2)$; the scalar factors of order $p^m$, $m\ge1$,
are analytic with derivative $O(p^{m-1})$. On the star at most two terms of size $O(h^3)$ are lost. The
products with $\mu$ use $\mu=O(h)$ and $\delta\mu=O(h^2)$; those of the chosen site use $\hat c=O(h)$ and
$\delta\hat c=O(h^2)$. In item~11, $\delta(6p^2+2p\alpha)=O(h^3)$; the interior sum uses
$\sum\tau|\mathsf G(x,e_0)|=O(1)$, and on the removed sites the mass is at most two and $\mathsf
G(x,e_0)=O(h)$.

For the trace we may keep $W$ exact, at a cost $O(h^3)$ on each board. Since $W$ annihilates the gauge,
the quadratic trace is $\operatorname{tr}(G^0_{II}WG^0_{II}W)$. The entries of $W$ in Step~7 give
$\max|W|=O(h^2)$, $\|W\|_\infty=O(h)$, and, on the common indices, $\max|\delta W|=O(h^3)$ and $\|\delta
W\|_\infty=O(h^2)$; here one uses $\delta s=O(h^2)$, $\delta p=O(ph)$ and the formula for $Q_J$. Write
$G^0_{II}=I+H_1$, with $\max|H_1|=O(h)$ and $\max|\delta H_1|=O(h^2)$. In $\operatorname{tr}(W^2)$,
changing $W$ costs $O(k\,h^2h^2)$ and losing a line costs $O(h\,h^2)$. The absolute sum in
$\operatorname{tr}(H_1W^2)$ is at most $k\max|H_1|\,\|W\|_\infty^2$, and that in
$\operatorname{tr}(H_1WH_1W)$ is at most $k^2\max|H_1|^2\|W\|_\infty^2$; changing one factor or
fixing a lost index improves each majorant by a factor $h$. Altogether $\delta r_{e_0}=O(h^3)$, which is
(a).
\end{proof}

\begin{corolario}[Increment of the two-cell remainder]\label{cor:incremento-marginal}
Under the hypotheses of Proposition~\ref{prop:menor-lejano}\textup{(b)}, let $R_\pi=\pi^T-\pi_1$ for the
pair $\{e,f\}$. If $c$ lies in none of the rows or columns of $e,f$, then
\[
 |R_\pi(M_c)-R_\pi(M)|\le C\,\frac{p_ep_f}{k^3}.
\]
\end{corolario}
\begin{proof}
By~\eqref{eq:psi-dos-casillas}, $R_\pi=p_ep_f\Psi$ with
$\Psi=(1+x_1)(1-g+t_1)(1+x_2+t_2)e^{\varrho_1+\varrho_2}-(1-g+x_1+x_2)$, where $x_1=\hat c_e$,
$x_2=\hat c_f$, $g=\mathsf G(e,f)$, $t_1=p'_f/p_f-1+g$, $t_2=\hat c'_f-\hat c_f$, $\varrho_1=r_e(M)$ and
$\varrho_2=r_f(M_e)$. Here $x_i,g=O(h)$ and $t_i,\varrho_i=O(h^2)$. The part of $\Psi$ linear in
$(x_1,x_2,g)$ vanishes, so its derivatives in $x_i,g$ are $O(h)$, while those in $t_i,\varrho_i$ are
$1+O(h)$. Let $\delta$ be the variation from $M$ to $M_c$. Then $\delta x_i,\delta g=O(h^2)$
by~\eqref{eq:variaciones-base}, and $\delta\varrho_1,\delta\varrho_2=O(h^3)$ by
Proposition~\ref{prop:menor-lejano}\textup{(a)}, applied to $M$ and to $M_e$. For $t_1,t_2$ we need
second differences. The child's Green matrix satisfies $G'-G^0_{II}=(I+H_1)W(I+H_2)+C_q$, where the
entries of $H_i$ are $O(h)$ with variations $O(h^2)$, and $C_q$ is built from $q=G^0_{IJ}v_J=O(h)$,
$\delta q=O(h^2)$, $v_I^\top q=O(h)$ and the divisor $2k-2$, so that $\delta C_q=O(h^3)$. With the bounds
for $W$ and $\delta W$ of Step~3 this gives $\delta(G'-G^0_{II})_{ij}=O(h^3)$. For
$X_x=(G^0a_x)^\top\eta^X-\varepsilon_x$ with $\eta^X_i=\sum_{y\in i}\tau_y\varepsilon_y$, we have
$\delta\Delta p=O(ph^2)$ from $\Delta p=\tau s+O(ps^2)$ and~\eqref{eq:masas-hijo}, hence
$\delta\Delta\varepsilon=O(h^3)$ and $\delta\Delta\eta^X_i=O(h^3)$; in the contraction with $G^0a_x$, whose
$\ell^1$ norm is $O(1)$ and whose entries on the four removed lines are $O(h)$, this gives
$\delta\Delta X=O(h^3)$ and $\delta t_2=O(h^3)$. Finally
$t_1=-(\varsigma_f/p_f)\mathsf G(e,f)+(\tau_f/p_f)a_f^\top\zeta+[u(t_fe^{s_f})-p_f-\tau_fs_f]/p_f$ has
variation $O(h^3)$ by Step~2. Therefore $\delta\Psi=O(h^3)$. Since $\Psi=O(h^2)$ and
$\delta(p_ep_f)=O(p_ep_fh)$, the claim follows.
\end{proof}

